\documentclass[11pt, a4paper]{article}

\usepackage[T1]{fontenc}
\usepackage[utf8]{inputenc}
\usepackage{tgpagella}
\usepackage[spanish, es-noshorthands]{babel}
\usepackage{amsmath, amssymb, bm}
\usepackage{tabularx}
\usepackage{booktabs}
\usepackage[most]{tcolorbox}
\usepackage[margin=2.5cm]{geometry}
\usepackage{ragged2e}
\usepackage{fancyhdr}

\pagestyle{fancy}
\fancyhf{}
\setlength{\headheight}{16pt}
\lhead{\textbf{UCB Sede Tarija} | Ing. Mecatrónica}
\chead{}
\rhead{IMT-121 Circuitos Electrónicos I | Readiness Checklist}
\lfoot{Prof: M.Sc. Ing. Bernardo Quiroga Turdera}
\rfoot{Página \thepage}
\renewcommand{\headrulewidth}{0.4pt}

\definecolor{ucbblue}{RGB}{0, 51, 102}

\begin{document}

\begin{tcolorbox}[colback=ucbblue!5!white, colframe=ucbblue, title=\textbf{Checklist Formativo de Progreso Docente (Readiness Gate)}]
    \centering \textbf{Decisión de Transición Hacia Simulación Computacional (.tran)}
\end{tcolorbox}

\noindent \textit{Instrucciones: Evalúe el desempeño macro de la cohorte durante la sesión de consolidación en base al Ejercicio de Transferencia Independiente. El paso al análisis transitorio simulado se condiciona a la robustez de este marco conceptual.}

\vspace{0.4cm}
\renewcommand{\arraystretch}{1.6}
\noindent
\begin{tabularx}{\textwidth}{>{\RaggedRight\arraybackslash}X >{\centering\arraybackslash}p{2cm} >{\centering\arraybackslash}p{2cm} >{\centering\arraybackslash}p{2cm}}
\toprule
\textbf{Indicador de Evidencia Conceptual} & \textbf{Sí} & \textbf{Parcial} & \textbf{No} \\ 
\midrule
Interpreta físicamente $\frac{dv_C}{dt}$ (Tasa de cambio) & $\square$ & $\square$ & $\square$ \\ \addlinespace
Aplica correctamente la ley constitutiva $i_C = C \frac{dv_C}{dt}$ & $\square$ & $\square$ & $\square$ \\ \addlinespace
Usa el principio de continuidad de voltaje adecuadamente & $\square$ & $\square$ & $\square$ \\ \addlinespace
Distingue claramente los instantes $0^-, 0^+$ y $\infty$ & $\square$ & $\square$ & $\square$ \\ \addlinespace
Encuentra el valor DC final abriendo el capacitor ($t \to \infty$) & $\square$ & $\square$ & $\square$ \\ \addlinespace
Determina $R_{eq}$ mirando estrictamente desde el puerto del capacitor & $\square$ & $\square$ & $\square$ \\ \addlinespace
Calcula $\tau$ de manera matemáticamente y dimensionalmente correcta & $\square$ & $\square$ & $\square$ \\ \addlinespace
Ensambla la fórmula de respuesta general exponencial & $\square$ & $\square$ & $\square$ \\ \addlinespace
Verifica matemáticamente los límites en $t=0^+$ y $t \to \infty$ & $\square$ & $\square$ & $\square$ \\ \addlinespace
Bosqueja la tendencia cualitativa correcta (Sube/Baja) & $\square$ & $\square$ & $\square$ \\ 
\bottomrule
\end{tabularx}

\vspace{0.8cm}
\begin{tcolorbox}[colback=gray!5, colframe=gray!50, title=\textbf{Fallo Final (Status de Transición)}]
Marque la directiva a ejecutar en la próxima sesión en base a la concentración de marcas (Checkmarks) de la tabla superior:
\begin{itemize}
    \item[$\square$] \textbf{\texttt{READY\_FOR\_TRANSIENT\_SIMULATION}}: Base fundacional sólida. Autorizado a enlazar analítica con LTSpice.
    \item[$\square$] \textbf{\texttt{NEEDS\_RC\_REINFORCEMENT}}: Deficiencias estructurales en KVL, continuidad o $R_{eq}$. Posponer la PC y repetir topologías RC guiadas.
    \item[$\square$] \textbf{\texttt{INDIVIDUAL\_FOLLOW\_UP}}: Mayoría competente. Intervenciones quirúrgicas en horario de oficina o laboratorio.
\end{itemize}
\end{tcolorbox}

\end{document}
